\documentclass{article}
\usepackage[T1]{fontenc}
\usepackage{lmodern}
\usepackage{graphicx}
\usepackage{amsmath,amssymb}
\usepackage[a4paper,margin=2cm]{geometry}
\usepackage[absolute,overlay]{textpos}
\setlength{\TPHorizModule}{1mm}
\setlength{\TPVertModule}{1mm}
\usepackage[indonesian]{babel}
\usepackage{hyperref}
\setlength{\emergencystretch}{2em}
% unit: Halaman 32

\begin{document}

% === Halaman 32 (python/native) ===

32

• Cipher ini  berhasil dipecahkan oleh Babbage dan Kasiski pada 

pertengahan Abad 19 (akan dijelaskan pada materi selanjutnya). 

• Kasiski menguraikan langkah-Langkah untuk menemukan panjang kunci 

(bukan huruf-huruf kunci kunci ). 

• Vigènere Cipher digunakan oleh Tentara Konfiderasi (Confederate Army) 

pada Perang Sipil Amerika (American Civil war). 

• Perang Sipil terjadi setelah Vigènere Cipher berhasil dipecahkan.

\end{document}
